% !TEX program = lualatex
\documentclass[a4paper,11pt]{ltjsarticle}

\usepackage{amsmath,amssymb,bm}
\usepackage{booktabs}
\usepackage{geometry}
\usepackage{graphicx}
\usepackage{xcolor}
\usepackage{hyperref}
\geometry{margin=24mm}

\title{\(\beta_{\mathrm{EI}}\) における EI の扱い\\
bo\_current との違いの整理}
\author{}
\date{2026年6月2日}

\newcommand{\BO}{\mathrm{BO}}
\newcommand{\GP}{\mathrm{GP}}
\newcommand{\EI}{\mathrm{EI}}
\newcommand{\HV}{\mathrm{HV}}
\newcommand{\Normal}{\mathcal{N}}

\begin{document}
\maketitle

\section{本メモの目的}

本メモは，今後の提案手法候補として扱う
\(\beta_{\mathrm{EI}}\) について，
BO 制御器内で用いる Expected Improvement（EI）の意味と，
現行版である \texttt{bo\_current} との違いを整理するものである．

本研究の提案手法では，多目的 GP の世代状態を観測し，
文脈付き BO により
\[
  u_\ell = (p_{c,\ell},p_{m,\ell},k_\ell)
\]
を決定する．
ここで \(p_c\) は交叉率，\(p_m\) は突然変異率，
\(k\) は次回の制御更新までの世代数である．

\section{BO 制御器の基本構造}

制御更新ステップ \(\ell\) において，
現在の GP 状態を文脈ベクトル
\[
  x_\ell =
  [
  \tau_\ell,
  \widetilde{\HV}_\ell,
  \widetilde{\Delta \HV}_\ell,
  D_\ell,
  \widetilde{L}_\ell,
  \widetilde{s}_\ell
  ]^\top
\]
として観測する．

BO 制御器は，この文脈 \(x_\ell\) の下で，
候補となる \((p_c,p_m,k)\) を評価し，
次に試す制御入力を選択する．
実装上は，\(k\) を連続変数として扱うのではなく，
\[
  k \in \mathcal{K}=\{1,3,5\}
\]
の離散候補として扱い，各 \(k\) ごとに別々の GP surrogate を持つ．
したがって，
\[
  f_k(x,p_c,p_m)
  \sim
  \GP
\]
を \(k=1,3,5\) それぞれについて学習する形である．

各 \(k\) に対して，現在文脈 \(x_\ell\) を固定し，
\((p_c,p_m)\) の候補集合を生成する．
その候補に対して GP surrogate は，
\[
  r \mid x_\ell,p_c,p_m,k
  \sim
  \Normal(\mu_k,\sigma_k^2)
\]
を予測する．
ここで \(r\) は，区間実行後に計算される報酬である．

図\ref{fig:beta-ei-control-flow}に，\texttt{bo\_current} と
\(\beta_{\mathrm{EI}}\) の違いを制御器内部の流れとして示す．
両者とも，現在状態 \(x_\ell\) を観測し，
\(k\in\{1,3,5\}\) の各候補について surrogate を用いて
EI を評価する点は同じである．
違うのは，EI を計算するときの基準値である．
\texttt{bo\_current} は \(k\) ごとの最良値を使い，
\(\beta_{\mathrm{EI}}\) は全 \(k\) 共通の最良値を使う．

\begin{figure}[htbp]
  \centering
  \includegraphics[width=\linewidth]{notes/figures/beta_EI/beta_EI_control_flow.png}
  \caption{\(\beta_{\mathrm{EI}}\) における EI 基準値変更の位置づけ}
  \label{fig:beta-ei-control-flow}
\end{figure}

\section{EI の基本式}

EI は Expected Improvement の略であり，
「ある候補を試したとき，過去の最良報酬をどの程度上回れそうか」
を表す獲得関数である．

ある候補点での予測平均を \(\mu\)，予測標準偏差を \(\sigma\)，
過去の基準値を \(r_{\mathrm{best}}\)，
探索余裕を \(\xi\) とすると，改善量は
\[
  I = \mu - r_{\mathrm{best}} - \xi
\]
である．
EI は
\[
  \EI = \mathbb{E}[\max(0,I)]
\]
として定義される．
GP surrogate の予測分布が正規分布であるとき，
\[
  z = \frac{\mu-r_{\mathrm{best}}-\xi}{\sigma}
\]
を用いて，
\[
  \EI
  =
  (\mu-r_{\mathrm{best}}-\xi)\Phi(z)
  +
  \sigma\phi(z)
\]
と計算できる．
ここで \(\Phi(\cdot)\) は標準正規分布の累積分布関数，
\(\phi(\cdot)\) は標準正規分布の確率密度関数である．

この式は，次の 2 つの効果を同時に含む．
\begin{itemize}
  \item 予測平均 \(\mu\) が高い候補を選びやすい．
  \item 予測不確実性 \(\sigma\) が大きい候補も，探索価値があるため選ばれやすい．
\end{itemize}

したがって EI は， exploitation（良さそうな候補の利用）と
exploration（不確実な候補の探索）のバランスを取る獲得関数である．

\section{\texttt{bo\_current} における EI}

\texttt{bo\_current} では，
EI の基準値 \(r_{\mathrm{best}}\) を
各 \(k\) の中での過去最高報酬として定義する．
すなわち，
\[
  r_{\mathrm{best}}^{(k)}
  =
  \max
  \{ r_i \mid k_i = k \}
\]
である．

この場合，\(k=1\) の候補を評価するときは，
\[
  r_{\mathrm{best}}^{(1)}
\]
を基準に EI を計算する．
同様に，\(k=3\) の候補では
\[
  r_{\mathrm{best}}^{(3)}
\]
を，\(k=5\) の候補では
\[
  r_{\mathrm{best}}^{(5)}
\]
を用いる．

したがって，\texttt{bo\_current} の EI は
\[
  \EI_k(x,p_c,p_m)
  =
  \EI
  \left(
  \mu_k(x,p_c,p_m),
  \sigma_k(x,p_c,p_m),
  r_{\mathrm{best}}^{(k)}
  \right)
\]
として解釈できる．

\subsection{直感的な意味}

\texttt{bo\_current} は，
各 \(k\) について，
\begin{quote}
  「この \(k\) の中で，今までより良くなりそうか」
\end{quote}
を見る方法である．

この設計の利点は，\(k=1,3,5\) をそれぞれ独立に育てられることである．
まだ全体としては弱い \(k\) であっても，
その \(k\) 内で改善余地があれば選ばれる可能性が残る．

一方で，弱点もある．
各 \(k\) が別々の基準値を持つため，
\[
  \text{「その \(k\) の中では良いが，全体ではそれほど良くない」}
\]
候補が高い EI を持つ場合がある．
つまり，\(k\) 間の比較基準が完全にはそろっていない．

\section{beta\_EI における EI}

\(\beta_{\mathrm{EI}}\) では，warm-up の順序は \texttt{bo\_current} と同じままで，
EI の基準値だけを変更する．
具体的には，各 \(k\) の中の過去最高ではなく，
全 \(k\) を含めた過去最高報酬を基準値とする．

すなわち，
\[
  r_{\mathrm{best}}^{\mathrm{global}}
  =
  \max_i r_i
\]
である．

このとき，任意の \(k\) に対して，
\[
  \EI_k(x,p_c,p_m)
  =
  \EI
  \left(
  \mu_k(x,p_c,p_m),
  \sigma_k(x,p_c,p_m),
  r_{\mathrm{best}}^{\mathrm{global}}
  \right)
\]
を計算する．

したがって，\(\beta_{\mathrm{EI}}\) は，
\begin{quote}
  「\(k\) が何であっても，全体の過去最高を超えられそうか」
\end{quote}
を見る方法である．

\section{両者の違い}

\begin{table}[htbp]
\centering
\caption{\texttt{bo\_current} と \(\beta_{\mathrm{EI}}\) の違い}
\begin{tabular}{lll}
\toprule
手法 & warm-up & EI の基準値 \\
\midrule
\texttt{bo\_current} & sequential & \(k\) ごとの過去最高 \(r_{\mathrm{best}}^{(k)}\) \\
\(\beta_{\mathrm{EI}}\) & sequential & 全 \(k\) 共通の過去最高 \(r_{\mathrm{best}}^{\mathrm{global}}\) \\
\bottomrule
\end{tabular}
\end{table}

図\ref{fig:ei-reference-comparison}は，この違いを数値例として示したものである．
左図の \texttt{bo\_current} では，\(k=1\) の候補は
\(k=1\) 内の最良値 \(0.30\) と比較される．
そのため，予測報酬が \(0.35\) 程度であれば，
\(k=1\) の中では改善候補として見えやすい．
右図の \(\beta_{\mathrm{EI}}\) では，全 \(k\) 共通の最良値 \(0.60\) が基準になる．
この場合，同じ \(0.35\) の候補は全体最良を超えられないため，
EI は小さくなりやすい．

\begin{figure}[htbp]
  \centering
  \includegraphics[width=\linewidth]{notes/figures/beta_EI/ei_reference_comparison.png}
  \caption{\texttt{bo\_current} と \(\beta_{\mathrm{EI}}\) における EI 基準値の違い}
  \label{fig:ei-reference-comparison}
\end{figure}

\texttt{bo\_current} は，
\[
  r_{\mathrm{best}}^{(k)}
  =
  \max \{ r_i \mid k_i=k \}
\]
を使う．
一方，\(\beta_{\mathrm{EI}}\) は，
\[
  r_{\mathrm{best}}^{\mathrm{global}}
  =
  \max \{ r_i \mid i=1,\dots,n \}
\]
を使う．

したがって，違いは
\[
  \text{「\(k\) ごとに改善を見るか」}
\]
と
\[
  \text{「全 \(k\) 共通で改善を見るか」}
\]
である．

\section{具体例}

例えば，過去の最高報酬が次のように得られていたとする．
\[
  r_{\mathrm{best}}^{(1)}=0.30,\quad
  r_{\mathrm{best}}^{(3)}=0.45,\quad
  r_{\mathrm{best}}^{(5)}=0.60.
\]

このとき，\texttt{bo\_current} で \(k=1\) の候補を評価する場合，
基準値は
\[
  r_{\mathrm{best}}^{(1)}=0.30
\]
である．
したがって，ある \(k=1\) 候補の予測報酬が \(0.35\) 程度であれば，
「\(k=1\) の中では改善が期待できる」と判断される．

一方，\(\beta_{\mathrm{EI}}\) では全体最高の
\[
  r_{\mathrm{best}}^{\mathrm{global}}=0.60
\]
が基準となる．
そのため，同じ \(k=1\) 候補が \(0.35\) 程度を予測されても，
全体最高の \(0.60\) は超えられないため，
EI は小さくなりやすい．

このように，\(\beta_{\mathrm{EI}}\) は，
各 \(k\) 内での局所的な改善ではなく，
全体として強い候補を優先しやすい．

\section{実験結果から見た位置づけ}

100 seed の比較では，Friedman-I において
\(\beta_{\mathrm{EI}}\) が平均 final archive HV で最良であった．

\begin{table}[htbp]
\centering
\caption{Friedman-I における 100 seed 比較}
\begin{tabular}{lrrrr}
\toprule
手法 & HV 平均 & HV 標準偏差 & Diversity 平均 & Diversity 標準偏差 \\
\midrule
\texttt{bo\_current} & 0.773081 & 0.055720 & 0.603050 & 0.123937 \\
\(\beta_k\) & 0.774463 & 0.046293 & 0.603639 & 0.119035 \\
\(\beta_{\mathrm{EI}}\) & 0.774922 & 0.050454 & 0.606697 & 0.122603 \\
\(\beta_{k\mathrm{EI}}\) & 0.771281 & 0.066523 & 0.600508 & 0.139885 \\
\bottomrule
\end{tabular}
\end{table}

図\ref{fig:seed100-metrics}に，Friedman-I と Poly-10 の
final archive HV および diversity を平均 \(\pm\) 標準偏差で示す．
Friedman-I では \(\beta_{\mathrm{EI}}\) が HV 平均で最も高く，
diversity もわずかに高い．
一方で，Poly-10 では \(\beta_{k\mathrm{EI}}\) の HV 平均が最も高いが，
各手法の差は標準偏差に比べると小さい．
このため，単純な平均値だけでなく，seed 内差分を見る必要がある．

\begin{figure}[htbp]
  \centering
  \includegraphics[width=\linewidth]{notes/figures/beta_EI/seed100_hv_diversity_bars.png}
  \caption{seed 100 比較における final archive HV と diversity}
  \label{fig:seed100-metrics}
\end{figure}

図\ref{fig:seed100-hv-diff}は，
同じ seed 内で \texttt{bo\_current} との差分を取ったものである．
Friedman-I では \(\beta_{\mathrm{EI}}\) の平均差が正であり，
かつ差分の標準偏差が他の変更案より小さい．
これは，\(\beta_{\mathrm{EI}}\) が少なくとも Friedman-I では
比較的安定した候補であることを示唆する．
ただし，差分そのものは小さいため，
今後は統計検定や \(k\) ごとの報酬分布の確認が必要である．

\begin{figure}[htbp]
  \centering
  \includegraphics[width=\linewidth]{notes/figures/beta_EI/seed100_hv_diff_vs_bo_current.png}
  \caption{\texttt{bo\_current} を基準にした seed 内 final HV 差分}
  \label{fig:seed100-hv-diff}
\end{figure}

この結果から，Friedman-I では，
\texttt{bo\_current} のように \(k\) ごとの局所的な改善を見るよりも，
\(\beta_{\mathrm{EI}}\) のように全 \(k\) 共通の基準で改善を評価した方が，
やや良い結果になった可能性がある．

ただし，差は非常に大きいわけではない．
\(\beta_{\mathrm{EI}}\) と \texttt{bo\_current} の HV 平均差は約 \(0.00184\) である．
したがって，今後は paired difference の可視化や統計検定により，
この差が偶然ではないかを確認する必要がある．

\section{beta\_EI の利点と注意点}

\subsection{利点}

\(\beta_{\mathrm{EI}}\) の利点は，\(k\) 間の EI 比較基準をそろえられることである．
\texttt{bo\_current} では，\(k=1,3,5\) がそれぞれ異なる過去最高を持つため，
各 \(k\) の EI が完全には同じ意味を持たない．
一方，\(\beta_{\mathrm{EI}}\) では全 \(k\) 共通の過去最高を使うため，
候補が
\[
  \text{「全体として本当に高い報酬を超えられそうか」}
\]
を基準に評価される．

これにより，特定の \(k\) 内だけでの局所的な改善に引っ張られにくくなる．

\subsection{注意点}

一方で，\(\beta_{\mathrm{EI}}\) にも注意点がある．
\(k=1,3,5\) は同じ報酬関数で評価されているが，
実際には観測ノイズ，更新頻度，区間内の多様性変化が異なる．
そのため，全 \(k\) 共通の過去最高をそのまま基準にすると，
ある \(k\) が不利になりすぎる可能性がある．

特に，ある \(k\) で偶然高い報酬が得られると，
その値が全体基準となり，他の \(k\) の候補の EI が小さく見えやすい．
したがって，\(\beta_{\mathrm{EI}}\) を本採用する場合でも，
次のような追加検討が必要である．

\begin{itemize}
  \item \(k\) ごとの報酬分布の違いを確認する．
  \item \(k\) ごとに標準化した EI を検討する．
  \item global EI と per-\(k\) EI の中間的な基準を検討する．
  \item paired difference や統計検定で，\(\beta_{\mathrm{EI}}\) の改善が安定しているか確認する．
\end{itemize}

\section{まとめ}

\texttt{bo\_current} と \(\beta_{\mathrm{EI}}\) の違いは，
EI の基準値をどこから取るかにある．

\[
\begin{aligned}
\texttt{bo\_current}:&
\quad
r_{\mathrm{best}}=r_{\mathrm{best}}^{(k)}
=
\max\{r_i\mid k_i=k\},
\\
\beta_{\mathrm{EI}}:&
\quad
r_{\mathrm{best}}=r_{\mathrm{best}}^{\mathrm{global}}
=
\max_i r_i.
\end{aligned}
\]

\texttt{bo\_current} は，各 \(k\) の中での改善を見やすい．
\(\beta_{\mathrm{EI}}\) は，全 \(k\) 共通の基準で候補を比較しやすい．

今回の seed 100 実験では，Friedman-I において
\(\beta_{\mathrm{EI}}\) が平均 HV で最も良かったため，
次段階の候補として優先的に検討する価値がある．
ただし，差は小さいため，
今後は seed 内差分，統計検定，\(k\) ごとの報酬分布分析を行い，
\(\beta_{\mathrm{EI}}\) の妥当性をより慎重に確認する必要がある．

\end{document}
